\documentclass{article}

% packages
\usepackage{amsmath, amsthm, thmtools, amsfonts, amssymb, luacode, catchfile, hyperref, ifthen, graphicx, bbm}
\ifcsname c@kobocompile\endcsname
	\usepackage[a5paper, total={1072pt, 1448pt}, margin=10pt, includeheadfoot]{geometry} % set page margins
\else
	\usepackage[a4paper, margin=50pt, includeheadfoot]{geometry}
\fi
\usepackage[shortlabels]{enumitem}
\usepackage[skip=3pt, indent=0pt]{parskip}

% language
\usepackage[bidi=basic, layout=tabular, provide=*]{babel}
\ifcsname c@english\endcsname
	\babelprovide[main, import]{english}
\else
	\babelprovide[main, import, alph=letters.gershayim]{hebrew}
	\setlist[enumerate,2]{label={\alph*.}}
	\babelprovide{rl}
\fi
\babelfont{rm}[Renderer=Harfbuzz]{Libertinus Serif}
\babelfont{sf}{Libertinus Sans}
\babelfont{tt}{Libertinus Mono}

% style
\ifcsname c@stoppagebreaks\endcsname
\else
	\AddToHook{cmd/section/before}{\clearpage}	% Add line break before section
\fi
\linespread{1.3}
\setcounter{secnumdepth}{0}		% Remove default number tags from sections, this won't do well with theorems
\AtBeginDocument{\setlength{\belowdisplayskip}{3pt}}
\AtBeginDocument{\setlength{\abovedisplayskip}{3pt}}
\graphicspath{ {../images/} }

% operators
\DeclareMathOperator\cis{cis}
\DeclareMathOperator\Sp{Span}
\DeclareMathOperator\tr{tr}
\DeclareMathOperator\im{Im}
\DeclareMathOperator\re{Re}
\DeclareMathOperator\diag{diag}
\DeclareMathOperator*\lowlim{\underline{lim}}
\DeclareMathOperator*\uplim{\overline{lim}}
\DeclareMathOperator\rng{rng}
\DeclareMathOperator\Sym{Sym}
\DeclareMathOperator\Arg{Arg}
\DeclareMathOperator\Log{Log}
\DeclareMathOperator\dom{dom}
\DeclareMathOperator\supp{Supp}
\DeclareMathOperator\var{Var}
\DeclareMathOperator\cov{Cov}
\DeclareMathOperator\End{End}
\DeclareMathOperator\id{id}
\DeclareMathOperator\cl{cl}
\DeclareMathOperator\gal{Gal}
\DeclareMathOperator\dist{dist}
\DeclareMathOperator\form{form}
\DeclareMathOperator\term{term}
\DeclareMathOperator\sent{sent}
\DeclareMathOperator\Th{Th}
\DeclareMathOperator\Hom{Hom}
\DeclareMathOperator\Aut{Aut}
\DeclareMathOperator\Char{char}
\DeclareMathOperator\vol{vol}
\DeclareMathOperator\GL{GL}
\DeclareMathOperator\SL{SL}
\DeclareMathOperator\lcm{lcm}

% commands
\newcommand{\cupdot}{\mathbin{\mathaccent\cdot\cup}}
\newcommand{\indicator}{\mathbbm{1}}
\newcommand{\AAA}[0]{\mathbb{A}}
\newcommand{\Aa}[0]{\mathcal{A}}
\newcommand{\BB}[0]{\mathbb{B}}
\newcommand{\Bb}[0]{\mathcal{B}}
\newcommand{\CC}[0]{\mathbb{C}}
\newcommand{\Cc}[0]{\mathcal{C}}
\newcommand{\DD}[0]{\mathbb{D}}
\newcommand{\Dd}[0]{\mathcal{D}}
\newcommand{\EE}[0]{\mathbb{E}}
\newcommand{\FF}[0]{\mathbb{F}}
\newcommand{\Ff}[0]{\mathcal{F}}
\newcommand{\Ii}[0]{\mathcal{I}}
\newcommand{\Gg}[0]{\mathcal{G}}
\newcommand{\HH}[0]{\mathbb{H}}
\newcommand{\Hh}[0]{\mathcal{H}}
\newcommand{\Kk}[0]{\mathcal{K}}
\newcommand{\Ll}[0]{\mathcal{L}}
\newcommand{\Mm}[0]{\mathcal{M}}
\newcommand{\NN}[0]{\mathbb{N}}
\newcommand{\Nn}[0]{\mathcal{N}}
\newcommand{\PP}[0]{\mathbb{P}}
\newcommand{\Pp}[0]{\mathcal{P}}
\newcommand{\QQ}[0]{\mathbb{Q}}
\newcommand{\RR}[0]{\mathbb{R}}
\newcommand{\Rr}[0]{\mathcal{R}}
\newcommand{\Ss}[0]{\mathcal{S}}
\newcommand{\TT}[0]{\mathbb{T}}
\newcommand{\Tt}[0]{\mathcal{T}}
\newcommand{\UU}[0]{\mathbb{U}}
\newcommand{\Uu}[0]{\mathcal{U}}
\newcommand{\Vv}[0]{\mathcal{V}}
\newcommand{\Ww}[0]{\mathcal{W}}
\newcommand{\Xx}[0]{\mathcal{X}}
\newcommand{\Yy}[0]{\mathcal{Y}}
\newcommand{\ZZ}[0]{\mathbb{Z}}
\newcommand{\Zz}[0]{\mathcal{Z}}
\newcommand{\acts}[0]{\circlearrowright}
\newcommand{\explain}[2] {
	\begin{flalign*}
		 && \text{#2} && \text{#1}
	\end{flalign*}
}
\newcommand{\maketitleprint}[1][green]{
	\begin{center}
		\includegraphics[width=6cm]{dragon_#1}
	\end{center}
}

% theorem commands
\newtheoremstyle{c_remark}
	{}	% Space above
	{}	% Space below
	{}% Body font
	{}	% Indent amount
	{\bfseries}	% Theorem head font
	{}	% Punctuation after theorem head
	{.5em}	% Space after theorem head
	{\thmname{#1}\thmnumber{ #2}\thmnote{ \normalfont{\text{(#3)}}}}	% head content
\newtheoremstyle{c_definition}
	{3pt}	% Space above
	{3pt}	% Space below
	{}% Body font
	{}	% Indent amount
	{\bfseries}	% Theorem head font
	{}	% Punctuation after theorem head
	{.5em}	% Space after theorem head
	{\thmname{#1}\thmnumber{ #2}\thmnote{ \normalfont{\text{(#3)}}}}	% head content
\newtheoremstyle{c_plain}
	{3pt}	% Space above
	{3pt}	% Space below
	{\itshape}% Body font
	{}	% Indent amount
	{\bfseries}	% Theorem head font
	{}	% Punctuation after theorem head
	{.5em}	% Space after theorem head
	{\thmname{#1}\thmnumber{ #2}\thmnote{ \text{(#3)}}}	% head content

\ifcsname c@english\endcsname
	\theoremstyle{plain}
	\newtheorem{theorem}{Theorem}[section]
	\newtheorem*{theorem*}{Theorem}
	\newtheorem{lemma}[theorem]{Lemma}
	\newtheorem{proposition}[theorem]{Proposition}
	\newtheorem*{proposition*}{Proposition}

	\theoremstyle{definition}
	\newtheorem{definition}[theorem]{Definition}
	\newtheorem*{definition*}{Definition}
	\newtheorem{example}{Example}[section]
	\newtheorem{exercise}{Exercise}[section]

	\theoremstyle{remark}
	\newtheorem*{remark}{Remark}
	\newtheorem*{solution}{Solution}
	\newtheorem{corollary}[theorem]{Corollary}
	\newtheorem{notation}[theorem]{Notation}
\else
	\theoremstyle{c_plain}
	\newtheorem{theorem}{משפט}[section]
	\newtheorem*{theorem*}{משפט}
	\newtheorem{lemma}[theorem]{למה}
	\newtheorem{proposition}[theorem]{טענה}
	\newtheorem*{proposition*}{טענה}

	\theoremstyle{c_definition}
	\newtheorem{definition}[theorem]{הגדרה}
	\newtheorem*{definition*}{הגדרה}
	\newtheorem{example}{דוגמה}[section]
	\newtheorem{exercise}{תרגיל}[section]

	\theoremstyle{c_remark}
	\newtheorem*{remark}{הערה}
	\newtheorem*{solution}{פתרון}
	\newtheorem{corollary}[theorem]{מסקנה}
	\newtheorem{notation}[theorem]{סימון}
\fi

% Questions related commands
\newcounter{question}
\setcounter{question}{1}
\newcounter{sub_question}
\setcounter{sub_question}{1}

\ifcsname c@english\endcsname
	\newcommand{\question}[1][0]{
		\ifthenelse{#1 = 0}{}{\setcounter{question}{#1}}
		\section{Question \arabic{question}}
		\addtocounter{question}{1}
		\setcounter{sub_question}{1}
	}

	\newcommand{\subquestion}[1][0]{
		\ifthenelse{#1 = 0}{}{\setcounter{sub_question}{#1}}
		\subsection{Part \alph{sub_question}}
		\addtocounter{sub_question}{1}
	}
\else
	\newcommand{\question}[1][0]{
		\ifthenelse{#1 = 0}{}{\setcounter{question}{#1}}
		\section{שאלה \arabic{question}}
		\addtocounter{question}{1}
		\setcounter{sub_question}{1}
	}

	\newcommand{\subquestion}[1][0]{
		\ifthenelse{#1 = 0}{}{\setcounter{sub_question}{#1}}
		\subsection{סעיף \localecounter{letters.gershayim}{sub_question}}
		\addtocounter{sub_question}{1}
	}
\fi

% import lua and start of document
\directlua{common = require ('../common')}

\GetEnv{AUTHOR}

% headers
\author{\AUTHOR}
\date\today

\title{יסודות באלגברה קומוטטיבית --- סיכום}
\setcounter{secnumdepth}{2}
% chktex-file 9
% chktex-file 17

\usepackage{fancyhdr}
\pagestyle{fancy}
\renewcommand{\headrulewidth}{0pt}

\begin{document}
\maketitle
\maketitleprint[red]{}

\tableofcontents

\section{שיעור 1 --- 18.10.2026}
אלגברה קומוטטיבית היא למעשה התורה של חוגים חילופיים (commutative).
היסטורית האלגברה הקומוטטיבית קיימת עבור תורת המספרים, הגאומטריה האלגברית ותורת האינווריאנטים אליה לא ניכנס.

\subsection{מבוא}
המושג הראשי שעליו נדבר בקורס הוא חוג חילופי (commutative ring), אובייקט הדומה לשדה אך ללא חילוק, ואידאל (ideal), המקבילה של תורת החוגים עבור תת־החבורה בתורת החבורות.
נעסוק גם במודול (module), המקבילה של מרחב וקטורי אך מעל חוג סקלרים במקום שדה סקלרים.

הקורס מבוסס בעיקר על הספר של Atiyah Macdoald שנקרא Introduction to commutative algebra, פרקים 1 עד 10.
\begin{definition}[חוג]
	הוא קבוצה עם שתי פעולות בינאריות, $(A, + : A \times A \to A, \cdot : A \times A \to A)$ המקיים את התכונות:
	\begin{enumerate}
		\item $(A, +)$ הוא חבורה חילופית, כלומר אסוציאטיבית, בעלת איבר נייטרלי (נסמן $0$) ולכל איבר יש איבר נגדי ביחס לפעולה
		\item חוק הקיבוץ בכפל: לכל $a, b, c \in A$ מתקיים $a (b c) = (a b) c$
		\item קיום איבר נייטרלי ביחס לכפל: קיים איבר $1 \in A$ כך ש־$1 \cdot a = a \cdot 1 = a$ לכל $a \in A$
		\item חלוק הפילוג (distributivity): לכל $a, b, c \in A$ מתקיים $(a + b) c = a c + b c$ וכן $a (b + c) = a b + a c$
	\end{enumerate}
\end{definition}
\begin{remark}
	האיבר הנייטרלי ביחס לכפל הוא יחיד, זאת שכן,
	\[
		1'
		= 1 \cdot 1'
		= 1
	.\]
\end{remark}
\begin{notation}
	ההופכי ביחס לחיבור יסומן $-a$, ובהתאם $a - b = a + (-b)$.
\end{notation}
\begin{definition}[חוג חילופי]
	חוג $(A, +, \cdot)$ נקרא חילופי (commutative) יש חילופיות גם ביחס לכפל, כלומר,
	\[
		\forall a, b \in A,\ 
		a \cdot b = b \cdot a
	.\]
\end{definition}
\begin{remark}
	חילופיות החיבור נובעת משאר התכונות.
	זאת שכן,
	\[
		(a + b) (1 + 1)
		= (a + b) 1 + (a + b) 1
		= a \cdot 1 + b \cdot 1 + a \cdot 1 + b \cdot 1
		= a + b + a + b
	.\]
	אבל מהצד השני,
	\[
		(1 + 1) (a + b)
		= a + a + b + b
	.\]
	$(A, +)$ היא חבורה ולכן נוכל לצמצם ולקבל בדיוק ש־$a + b = b + a$.
\end{remark}
\begin{remark}
	מקיבוציות הכפל נובע שמכפלה $a_1 \cdot a_k$ אינה תלויה בהצבת סוגריים, ולכן מותר להשמיט את הסוגריים.
\end{remark}
\begin{example}
	חוג השלמים $\ZZ$ עם פעולות הכפל והחיבור הסטנדרטיות הוא חוג חילופי.
\end{example}
\begin{example}
	לכל שדה $\FF$, $(\FF, +, \cdot)$ הוא חוג חילופי גם כן.
\end{example}
\begin{example}
	חוג המטריצות $A = M_n(\FF)$, כלומר המטריצות מעל $\FF$ מגודל $n \times n$ כאשר $1_A = I_n$ ו־$0_A = 0$.
	זהו חוג שחילופי אם ורק אם $n = 1$.
	נבחין גם כי במקום שדה $\FF$ יכולנו לבחור חוג כללי.
\end{example}
\begin{example}
	חוג הפולינומים,
	יהי $\FF$ שדה, נסמן את $\FF[x]$ להיות אוסף הפולינומים במשתנה $x$ עם מקדמים ב־$\FF$, וזהו חוג חילופי.
\end{example}
\begin{example}
	החוג הטריוויאלי:
	יהי $A$ יחידון, אז נגדיר את $+, \cdot$ טריוויאלית ונקבל ש־$(A, +, \cdot)$ הוא חוג חילופי.
	בחוג זה אף מתקיים $1_A = 0_A$, וזוהי תכונה המאפיינת ביחידות חוג זה, כלומר אם $0 = 1$ אז מתקיים,
	\[
		a
		= a (1 + 0)
		= a ( 1 + 1)
		= a + a
	.\]
	כלומר $a = 0$ לכל $a \in A$.
\end{example}

\subsection{פעולות בסיסיות בחוגים}
יהי $(A, +, \cdot)$ חוג.
\begin{definition}[חזקה]
	אם $n \in \NN$ ו־$a \in A$ אז נגדיר,
	\[
		a^n
		= \overbrace{a \cdots a}^{n \text{ times}}
	.\]
\end{definition}
\begin{proposition}[תכונות חזקה]
	לכל $n, m \in \NN$ ו־$a, b \in A$ מתקיים,
	\begin{enumerate}
		\item $a^{n + m} = a^n a^m$
		\item אם $A$ חילופי אז גם,
			\[
				{(a b)}^n = a^n b^n
			.\]
	\end{enumerate}
\end{proposition}
\begin{definition}[הכפלה במספר שלם]
	אם $n \in \NN$ ו־$a \in A$ אז נגדיר את,
	\[
		n a
		= \begin{cases}
			\overbrace{a + \cdots + a}^{n \text{times}} & n > 0 \\
			0 & n = 0 \\
			\overbrace{(-a) + \cdots + (-a)}^{n \text{times}} & n < 0
		\end{cases}
	.\]
\end{definition}
לכפל בשלם יש אסוציאטיביות.

\subsection{יחידות}
\begin{definition}[איבר הפיך]
	$a \in A$ נקרא הפיך (invertible) או יחידה (unit) אם קיים $b \in A$ כך שמתקיים $a b = b a = 1$.
	במקרה זה נאמר ש־$b$ הוא ההופכי של $a$ ונסמן $b = a^{-1}$.
\end{definition}
\begin{notation}
	נסמן את קבוצת ההפיכים על־ידי,
	\[
		A^{\times}
		= \{ a \in A \mid a \text{ is invertible} \}
	.\]
	סימון מקובל נוסף הוא $A^*$.
\end{notation}
\begin{proposition}
	$A^\times$ חבורה ביחס לכפל.
\end{proposition}
\begin{proof}
	נבדוק ש־$A^\times$ סגור ביחס לכפל.
	אם $a, b \in A^\times$ אז קיימים $a', b' \in A^\times$ המקיימים $a a' = b b' = 1$.
	מתקיים,
	\[
		(a b) (b' a')
		= a b b' a'
		= a a'
		= 1
	.\]
	וכן באופן דומה נקבל זאת עבור כפל הפוך, לכן $a b \in A^\times$.
	בהתאם נובע שכפל הוא פעולה בינארית על $A^\times$, וכן ברור כי $1 \in A^\times$ נייטרלי וכן שלכל $a \in A^\times$ גם $a' \in A^\times$.
\end{proof}
לחבורה $A^\times$ קוראים גם החבורה הכפלית של החוג (multiplication group).


\listoftheorems[title=הגדרות ומשפטים,ignoreall,show={theorem,definition},swapnumber,onlynamed={proposition}]
\end{document}
